\documentclass[11pt]{article}
\usepackage[a4paper,margin=1in]{geometry}
\usepackage[T1]{fontenc}
\usepackage{graphicx,amsmath,amssymb,booktabs,array,float}
\usepackage{xcolor,listings,caption,fancyhdr,needspace}
\usepackage[most]{tcolorbox}
\usepackage[hidelinks]{hyperref}
\usepackage{pgfplots}

\pgfplotsset{compat=1.18}
\definecolor{accent}{HTML}{31566B}
\definecolor{muted}{HTML}{65717A}
\definecolor{pale}{HTML}{F5F7F8}
\definecolor{border}{HTML}{DCE2E6}
\setcounter{secnumdepth}{0}
\setlength{\parindent}{0pt}
\setlength{\parskip}{0.4em}
\setlength{\emergencystretch}{2em}
\captionsetup{font=small,labelfont=bf,skip=7pt}
\renewcommand{\arraystretch}{1.55}
\pagestyle{fancy}
\fancyhf{}
\fancyhead[L]{\small EE 3150: Signals and Systems}
\fancyhead[R]{\small Lab 3}
\fancyfoot[C]{\thepage}
\setlength{\headheight}{15pt}
\renewcommand{\headrulewidth}{0.3pt}


\newcommand{\answer}[1][0.45cm]{%
  \begin{tcolorbox}[colback=pale,colframe=border,boxrule=0.4pt,
    arc=1mm,left=9pt,right=9pt,top=6pt,bottom=6pt,breakable]
  {\color{muted}\itshape [write here]}\par\vspace{#1}
  \end{tcolorbox}}
\newcommand{\blank}{\rule{1.65cm}{0.35pt}}
\newcommand{\topic}[1]{\Needspace{4\baselineskip}\par\textbf{#1}\par\nopagebreak}

\newcommand{\labfigure}[2]{%
  \begin{figure}[H]
    \centering
    \IfFileExists{#1}{\includegraphics[width=0.88\linewidth,height=0.65\textheight,keepaspectratio]{#1}}{%
    \IfFileExists{figs/#1}{\includegraphics[width=0.88\linewidth,height=0.65\textheight,keepaspectratio]{figs/#1}}{%
    \fcolorbox{border}{pale}{\parbox[c][1.1in][c]{0.87\linewidth}{%
      \centering\color{muted}\small Figure placeholder\\[6pt]
      \texttt{\detokenize{#1}}}}}}
    \caption{#2}
  \end{figure}}

\lstdefinestyle{python}{
  language=Python,basicstyle=\ttfamily\footnotesize,
  keywordstyle=\color{accent}\bfseries,
  commentstyle=\color{muted},stringstyle=\color{accent},
  backgroundcolor=\color{pale},frame=single,rulecolor=\color{border},
  breaklines=true,showstringspaces=false,tabsize=4,
  columns=fullflexible,keepspaces=true,
  xleftmargin=6pt,xrightmargin=6pt,framexleftmargin=6pt,
  framexrightmargin=6pt,aboveskip=8pt,belowskip=8pt,
  numbers=none,captionpos=b
}

\begin{document}

\begin{titlepage}
  \centering
  \vspace*{3cm}
  {\large EE 3150: Signals and Systems}\\[8pt]
  {\LARGE\bfseries Convolution and Filtering}\\[12pt]
  Colden Briggs, Eben Radocchia, and Liam Donlan\\[4pt]
  Instructor: Dr. Ossareh \quad $\vert$ \quad University of Vermont\\[4pt]
  \today
\end{titlepage}

\tableofcontents

\newpage


\section{Introduction}

The purpose of this lab was to understand convolution and
how it is used to filter signals. Discrete-time convolution
was calculated by hand and in Python. The step and cosine
responses of an RC circuit were then found, and the
two-stage response was compared with measurements from
Lab 2. Moving average, difference, and echo filters were
applied to an audio clip. Finally, an IIR filter was
implemented using a difference equation and compared with
the moving average and analytical RC response.

\newpage

\section{Part 1: Discrete-Time Convolution}
\subsection*{Task 1: Mechanics of Convolution}
The input sequences are
\[
  x[n]=\{1,2,3,1\},\quad n=0,\ldots,3,
  \qquad h[n]=\{1,-1,2\},\quad n=0,\ldots,2.
\]

\topic{Hand calculation}
\labfigure{task1_handwritten.pdf}{Handwritten flip-and-slide calculation of $x[n]*h[n]$.}
The convolution gives
\[
  y[n]=\{1,1,3,2,5,2\},\qquad n=0,\ldots,5.
\]
The expected output length is
\[
  N_y=N_x+N_h-1=4+3-1=6.
\]
The result has six samples, which matches the expected length.

\clearpage
\topic{Python code}
\begin{lstlisting}[style=python]
import numpy as np
import matplotlib.pyplot as plt

x = np.array([1, 2, 3, 1])
h = np.array([1, -1, 2])


# This code for convolution only works if x and h both start at time n=0
def conv_loop(x, h):
    Ny = len(x) + len(h)- 1
    y = np.zeros(Ny)
    for n in range(Ny): # this loop computes y[n] for different n's
        for k in range(len(x)): # this loop implements the convolution sum
            m = n - k
            if 0 <= m < len(h):
                y[n] = y[n] + x[k]*h[m]
    return y

#D.T. convolution output
y1 = conv_loop(x, h)
n1 = np.arange(len(y1))

#C.T. convolution output
y2 = np.convolve(x,h)
n2 = np.arange(len(y2))


#Plot results
fig1, ax1 = plt.subplots(figsize=(8, 6))
fig2, ax2 = plt.subplots(figsize=(8, 6))

ax1.stem(n1, y1, label="y1[n] = x[n] * h[n]")

ax1.set_title("Convolution output")
ax1.set_xlabel("n")
ax1.legend()
ax1.grid(True)
fig1.tight_layout()

ax2.stem(n2, y2, label="y2[n] = x[n] * h[n]")
ax2.set_title("Convolution output 2")
ax2.set_xlabel("n")
ax2.legend()
ax2.grid(True)
fig2.tight_layout()
plt.show()
\end{lstlisting}

\clearpage
\topic{Convolution results}
\labfigure{task1_conv.png}{$y[n]$ computed with \texttt{conv\_loop}.}
\labfigure{task1_conv2.png}{$y[n]$ computed with \texttt{np.convolve}.}
The handwritten result, the \texttt{conv\_loop} output, and the
\texttt{np.convolve} output all give $y[n]=\{1,1,3,2,5,2\}$.
The two stem plots are identical, so all three methods agree.

\topic{Convolution with a shifted impulse}
\[
  x[n]*\delta[n-2]=x[n-2]=\{0,0,1,2,3,1\},\qquad n=0,\ldots,5.
\]
\labfigure{task1_shifted_handwritten.png}{Hand calculation and sketch of $x[n]*\delta[n-2]$.}
Since $\delta[n-2]$ is nonzero only at $n=2$, the convolution sum
$\sum_k x[k]\delta[n-k-2]$ keeps only the term $k=n-2$, giving
$x[n-2]$, which is $x[n]$ delayed by two samples.

\clearpage
\section{Part 2: Continuous-Time Convolution}
\subsection{Task 2.1: Step Response of the First Stage}
\[
  h(t)=\frac{1}{\tau}e^{-t/\tau}u(t),\qquad \tau=R_1C_1=0.01\ \text{s}.
\]

\topic{Hand derivation and final response}
\labfigure{task2.1_hand_derivation.jpeg}{Hand derivation of the first-stage step response $u(t)*h(t)$.}

\topic{Numerical response}

\begin{lstlisting}[style=python]
import numpy as np
import matplotlib.pyplot as plt

tau = 0.01
dt = 0.0001
t = np.arange(0.0, 0.08, dt)

#impulse response 
h = (1/tau) * np.exp(-t/tau)

#unit step input
x = np.ones(len(t))

#output convolution of impulse respones and unit step
y = np.convolve(x, h) * dt

# use seconds for the output time axis
t_axis = np.arange(len(y)) * dt

#analytic step response
y_analytic = 1 - np.exp(-t/tau)

#plot and compare analytic and numerical convolution results
fig1, ax1 = plt.subplots(figsize=(8, 5))
ax1.plot(t, y_analytic, label="analytic step response")
ax1.plot(t_axis, y, "--", label="numeric step response")
ax1.set_xlim(0, 0.08)
ax1.set_ylim(0, 1.1)
ax1.set_title("First-stage RC step response")
ax1.set_xlabel("time (s)")
ax1.set_ylabel("output amplitude")
ax1.legend()
ax1.grid(True)
fig1.tight_layout()
plt.show()
\end{lstlisting}
\labfigure{task2.1_conv.png}{Numerical step response of the first RC stage.}
The hand-calculated analytic step response and the numerically calculated step response match each other very closely.

\topic{Filtering behavior}

The input jumps immediately from 0 to 1, but the output
rises gradually toward 1. The capacitor takes time to charge,
so the output cannot follow the sudden change immediately.
This smooths the input and shows the low-pass behavior of
the RC circuit. Slow changes pass through more easily,
while fast changes are reduced.


\clearpage
\subsection{Task 2.2: Cosine Response of the First Stage}
\[
  x(t)=\cos(\omega_0t)u(t),\qquad \tau=0.01\ \text{s}.
\]

\topic{Hand derivation}
\labfigure{task2.2_hand_derivations.png}{Hand derivation of the first-stage response to a cosine input.}
The required form of the response is
\[
 y(t)=\frac{\cos\!\left(\omega_0t-\tan^{-1}(\omega_0\tau)\right)}
 {\sqrt{1+(\omega_0\tau)^2}}
 -\frac{e^{-t/\tau}}{1+(\omega_0\tau)^2},\qquad t\geq 0.
\]

\topic{Steady-state response}
As $t \to \infty$, the exponential term goes to 0, leaving the steady state response as follows:
\[
 y(t)=\frac{\cos\!\left(\omega_0t-\tan^{-1}(\omega_0\tau)\right)}
 {\sqrt{1+(\omega_0\tau)^2}},\qquad t\geq 0.
\]

This is a cosine signal, with amplitude $\frac{1}{\sqrt{1+(\omega_0\tau)^2}}$, frequency $\omega_0$, and phase $-\tan^{-1}(\omega_0\tau)$

\begin{table}[H]
\centering
\begin{tabular}{ccc}
\toprule
$\omega_0\tau$ & Amplitude & Phase ($^\circ$)\\
0.1 & 0.99504 & $-5.711$\\
1   & 0.70711 & $-45.000$\\
10  & 0.09950 & $-84.289$\\
\bottomrule
\end{tabular}
\caption{Amplitude and phase of the first-stage output cosine.}
\end{table}

\topic{Effect of frequency}

At low frequencies, the output keeps most of the input
amplitude and has only a small phase shift. As frequency
increases, the output amplitude decreases and the phase
approaches $-90^\circ$. This shows that the circuit passes
low-frequency signals more strongly than high-frequency
signals, making it a low-pass filter.

At high frequencies, where $\omega_0\tau\gg1$, the amplitude
is approximately $1/(\omega_0\tau)$ and the phase is
approximately $-90^\circ$. The steady-state output is then
\[
y_{\mathrm{ss}}(t)
\approx
\frac{1}{\omega_0\tau}
\cos\left(\omega_0t-\frac{\pi}{2}\right)
=
\frac{1}{\omega_0\tau}\sin(\omega_0t).
\]

This has the form of the input's integral multiplied by
$1/\tau$:
\[
\frac{1}{\tau}\int\cos(\omega_0t)\,dt
=
\frac{1}{\omega_0\tau}\sin(\omega_0t).
\]
The circuit therefore acts approximately like a scaled
integrator at high frequencies.

\clearpage
\clearpage
\subsection{Task 2.3: The Two-Stage Cascade}

The two-stage circuit has
\[
\tau_1=0.01\ \text{s},
\qquad
\tau_2=0.0047\ \text{s},
\qquad
G=5.7,
\qquad
\omega=2\pi f.
\]


\topic{Cascade derivation}

From Task 2.2, the first RC stage changes the amplitude
and phase of a cosine input by
\[
A_1=\frac{1}{\sqrt{1+(\omega\tau_1)^2}},
\qquad
\phi_1=-\tan^{-1}(\omega\tau_1).
\]
For a unit-amplitude input, its steady-state output is
\[
v_{c1,\mathrm{ss}}(t)
=A_1\cos(\omega t+\phi_1).
\]

The gain buffer multiplies this signal by $G$ without
changing its phase:
\[
v_{\mathrm{buf,ss}}(t)
=GA_1\cos(\omega t+\phi_1).
\]

Applying the same result to the second RC stage gives
another amplitude factor and phase shift:
\[
A_2=\frac{1}{\sqrt{1+(\omega\tau_2)^2}},
\qquad
\phi_2=-\tan^{-1}(\omega\tau_2).
\]
The final steady-state output is therefore
\[
v_{c2,\mathrm{ss}}(t)
=GA_1A_2\cos(\omega t+\phi_1+\phi_2).
\]

The amplitude factors multiply, while the phase shifts add.
Substituting the expressions for $A_1$, $A_2$, $\phi_1$,
and $\phi_2$ gives the magnitude and phase below.
The resulting magnitude and phase are
\begin{align*}
M(\omega)
&=\frac{G}
{\sqrt{1+(\omega\tau_1)^2}
 \sqrt{1+(\omega\tau_2)^2}},\\
\phi(\omega)
&=-\tan^{-1}(\omega\tau_1)
  -\tan^{-1}(\omega\tau_2).
\end{align*}

\topic{Experimental and analytical magnitudes}

The experimental values were taken from the second-stage
output measurements in Task 3.4 of Lab 2. The input was
1 V peak-to-peak, so the experimental magnitude was calculated as
\[
M_{\mathrm{exp}}
=\frac{V_{\mathrm{out,pp}}}{V_{\mathrm{in,pp}}}
=\frac{V_{\mathrm{out,pp}}}{1\ \text{V}}.
\]

The analytical values were calculated using the magnitude
expression above.

\begin{table}[H]
\centering
\begin{tabular}{ccc}
\toprule
Frequency (Hz) & Experimental magnitude & Analytical magnitude\\
\midrule
5   & 5.37 & 5.380\\
15  & 3.80 & 3.793\\
30  & 2.01 & 1.999\\
60  & 0.71 & 0.718\\
120 & 0.20 & 0.203\\
\bottomrule
\end{tabular}
\caption{Experimental and analytical voltage gains of the
two-stage RC circuit.}
\end{table}

\labfigure{p2_two_stage_mag.png}{Experimental and analytical magnitude versus
frequency for the two-stage RC circuit.}


\topic{Experimental and analytical agreement}

The experimental and analytical values agree closely at
all five frequencies. For example, the measured magnitude
at 5 Hz was 5.37, compared with an analytical value of 5.380.
At 120 Hz, the measured magnitude was 0.20, compared with
0.203 analytically. Small differences could come from
component tolerances and measurement uncertainty.

\topic{Two-stage filtering behavior}

The magnitude decreases as frequency increases. Low-frequency
signals pass through with amplification, while higher-frequency
signals are strongly reduced. The measured magnitude falls
from 5.37 at 5 Hz to 0.20 at 120 Hz.

Both RC stages reduce high-frequency content, so their
combined response falls faster than the response of either
stage alone. The circuit therefore acts as a two-stage
low-pass filter.
\clearpage
\section{Part 3: FIR Audio Filters}

\subsection{Task 3.1: Audio Filtering}

The original audio had a sampling rate of $f_s=16000$ Hz.
Each filtered output was normalized before being saved and played.

\topic{Python code}
\begin{lstlisting}[style=python]
from scipy.io import wavfile
import numpy as np
import matplotlib.pyplot as plt

#read sampling frequency and data from audio clip
fs, clip = wavfile.read("clip.wav")
clip = clip.astype(float); 
if clip.ndim > 1: 
    clip = clip[:,0]
clip = clip/np.max(np.abs(clip))

#Moving average filter
h_ma = np.ones(64)/64

#Difference filter
h_diff = np.array([1.0, -1.0])


#Echo filter
D = int(round(0.2 * fs))
h_echo = np.zeros(D+1)
h_echo[0] = 1.0
h_echo[D] = 0.6

#Moving average filter output
y_ma = np.convolve(clip,h_ma)
y_ma = y_ma / np.max(np.abs(y_ma)) # normalize to [-1, 1]

#Difference filter output
y_diff = np.convolve(clip,h_diff)
y_diff = y_diff / np.max(np.abs(y_diff)) # normalize to [-1, 1]

#Echo filter output
y_echo = np.convolve(clip,h_echo)
y_echo = y_echo / np.max(np.abs(y_echo)) # normalize to [-1, 1]


#write outputs to files for comparison
from scipy.io import wavfile
y_ma_i16 = np.int16(y_ma / np.max(np.abs(y_ma)) * 32767) # converts to 16 bit int for wav format
fname = "low_pass.wav"
wavfile.write(fname, round(fs), y_ma_i16) # writes the wav file to disk

y_diff_i16 = np.int16(y_diff / np.max(np.abs(y_diff)) * 32767) # converts to 16 bit int for wav format
fname = "high_pass.wav"
wavfile.write(fname, round(fs), y_diff_i16) # writes the wav file to disk

y_echo_i16 = np.int16(y_echo / np.max(np.abs(y_echo)) * 32767) # converts to 16 bit int for wav format
fname = "echo.wav"
wavfile.write(fname, round(fs), y_echo_i16) # writes the wav file to disk

\end{lstlisting}

\topic{Moving average filter}

The moving average combines 64 samples. Fast changes
are averaged out, reducing high-frequency detail and giving
the audio a smoother, more muffled sound. It is called a
low-pass filter because lower frequencies pass through more
strongly than higher frequencies.

\topic{Difference filter}

The difference filter subtracts the previous sample from the
current sample. Slowly changing parts mostly cancel, while
rapid changes remain. This reduces low-frequency content and
emphasizes sharper changes, so the output is sounds
thinner and less bass-heavy. It is called a high-pass filter
because it emphasizes higher frequencies relative to lower
frequencies.

\topic{Echo filter}

For an echo delay of $0.2$ s, the number of samples is
\[
D=\operatorname{round}(0.2f_s)
 =\operatorname{round}(0.2\times16000)
 =3200.
\]

The first nonzero value of the impulse response passes the
original signal. The second produces a copy delayed by
3200 samples and scaled by $0.6$. The output therefore
contains the original sound followed by a quieter copy
$0.2$ s later.


\subsection{Task 3.2: Order Does Not Matter}

\topic{Python code}

The following code uses the original clip and filters
defined in Task 3.1.

\begin{lstlisting}[style=python]
# moving average followed by echo
y1 = np.convolve(np.convolve(clip, h_ma), h_echo)

# echo followed by moving average
y2 = np.convolve(np.convolve(clip, h_echo), h_ma)

# one combined filter
h_combined = np.convolve(h_ma, h_echo)
y3 = np.convolve(clip, h_combined)

# compare the outputs before normalization
print(np.allclose(y1, y2))
print(np.allclose(y1, y3))
\end{lstlisting}

\topic{Numerical agreement and sound}

Both comparisons return \texttt{True}. The largest difference
between the unnormalized outputs is approximately
$5.6\times10^{-17}$, which is just numerical rounding.
The saved audio files differ by at most one count in the
16-bit samples.

These results show that switching the order of the filters
or combining them first gives the same output to numerical
precision.

\topic{Meaning of the combined impulse response}

The expression \texttt{np.convolve(h\_ma, h\_echo)} gives the
impulse response of one filter that performs both operations.
Convolving the clip with this combined filter produces the
same result as passing it through the moving average and
echo filters separately.


\clearpage
\subsection{Task 3.3: Moving Average Gain versus Frequency}

One-second cosine signals were generated at the eight test
frequencies. Each signal had an input amplitude of 1 and was
filtered using the 64-sample moving average.

The output amplitude was measured using the maximum absolute
value in the second half of the full convolution output,
following the lab procedure.

\topic{Python code}
\begin{lstlisting}[style=python]
from scipy.io import wavfile
import numpy as np
import matplotlib.pyplot as plt

#sampling rate
fs = 16000

#test frequencies
freq = [50, 100, 200, 500, 1000, 2000, 4000, 8000]

t = np.arange(0, 1.0, 1.0/fs)

#create list of one second pure tones at each test frequency
tones = []
for f in freq:
    tones.append(np.cos(2*np.pi*f*t))


h_ma = np.ones(64)/64 #length-64 moving average filter      

#convolve each tone with length-64 moving average filter
y = []
for t in tones:
    y.append(np.convolve(t, h_ma))

#measure steady state amplitude of output at each frequency
steady_state_amp = []
for out in y:
    steady_state_amp.append(np.max(np.abs(out[len(out)//2:])))

#assign amplitude of lowest frequency to variable
amp_low_freq = steady_state_amp[0] 

# code for plotting freq vs amplitude
fig, ax = plt.subplots()

ax.plot(freq, steady_state_amp, 'o-')

ax.set_title("Amplitude vs Frequency (64-pt moving average filter output)")
    
ax.grid(True, alpha=0.3)
ax.set_xlabel("Frequency (Hz)")
ax.set_ylabel("Amplitude")
ax.set_xscale('log')

#Mark low frequency amplitude as well as half of low frequency amplitude
ax.axhline(y=amp_low_freq, color="tab:red", linestyle="--", linewidth=1, label="Low frequncy amplitude")
ax.axhline(y=amp_low_freq/2, color="tab:green", linestyle="--", linewidth=1, label="Half of low frequency amplitude")
ax.legend() 

plt.tight_layout()
fig.savefig("task3.3_amplitude_vs_freq.png", dpi=200)
plt.show()

\end{lstlisting}

\labfigure{task3.3_amplitude_vs_freq.png}
{Moving average output amplitude versus frequency on a
logarithmic frequency axis.}

\topic{Half-amplitude frequency}

The output amplitude at the lowest test frequency, 50 Hz,
was approximately $0.935$. Half of this value is
\[
\frac{0.935}{2}\approx0.468.
\]

The measured amplitude was approximately $0.757$ at 100 Hz
and $0.234$ at 200 Hz. The half-amplitude crossing is therefore
between these frequencies, at roughly
\[
f_{\mathrm{half}}\approx150\ \text{Hz}.
\]
This is an estimate because only eight frequencies were tested.

\topic{Agreement with the audio filtering}

The decreasing output amplitude is consistent with the
loss of high-frequency detail described in Task 3.1.
Reducing these frequencies makes the filtered audio
smoother and more muffled.

\topic{Low-pass behavior}

The output amplitude is larger at low frequencies and much
smaller at high frequencies. This shows that the moving
average acts as a low-pass filter.

The small nonzero values at some higher frequencies also
include the ending transient, since the specified amplitude
measurement includes the tail of the full convolution.
\clearpage
\clearpage
\section{Part 4: IIR Filtering}

The IIR filter computes each output sample using the previous
output and the current input:
\[
y[n]=a\,y[n-1]+b\,x[n],
\]
where
\[
a=\frac{\tau}{\tau+dt},
\qquad
b=\frac{dt}{\tau+dt},
\qquad
dt=\frac{1}{f_s}.
\]

\subsection{Task 4.1: IIR Streaming}

The step response was calculated using
\[
\tau=0.01\ \text{s},
\qquad
f_s=16000\ \text{Hz},
\qquad
dt=0.0000625\ \text{s}.
\]
These values give
\[
a\approx0.993789,
\qquad
b\approx0.006211.
\]

\topic{Python code}
\begin{lstlisting}[style=python]

import numpy as np
import matplotlib.pyplot as plt

#RC time constant
tau = 0.01

#sampling rate
fs = 16000

dt = 1.0/fs
a = tau/(tau+dt)
b = dt/(tau+dt)

x = np.ones(fs)
y = np.zeros_like(x)

#numeric step result convolution via DE approximation
for n in range(len(x)): # y[n] is the current sample of the live stream
    y[n] = a*y[n-1] + b*x[n]

fig, ax = plt.subplots()
n = np.arange(len(y))

#analytic step response
s_analytic = 1 - np.exp(-n*dt/tau)

#Plot and compare analytic and numeric step results
ax.plot(n, y, label="numeric")
ax.plot(n, s_analytic, "--", label="analytic")
ax.legend()

ax.set_title("Analytic vs. Numeric RC Step Response")

ax.grid(True, alpha=0.3)
ax.set_xlabel("Sample (n)")
ax.set_ylabel("Output")

plt.tight_layout()
fig.savefig("task4.1_step_response.png", dpi=200)
plt.show()
\end{lstlisting}

\labfigure{task4.1_step_response.png}
{Numerical IIR step response compared with the analytical
RC step response.}

\topic{Comparison of the step responses}

The numerical and analytical curves nearly overlap. Both
rise quickly at first and then approach a final value of 1.
The small difference comes from approximating the
continuous-time system using a finite sampling interval.
This shows that the IIR filter closely approximates the
RC circuit's step response.


\clearpage
\subsection{Task 4.2: Audio Filtering and FIR/IIR Comparison}

The audio filter used
\[
R_1=3.3\ \text{k}\Omega,
\qquad
C_1=470\ \text{nF},
\qquad
\tau\approx1.55\ \text{ms}.
\]

At $f_s=16000$ Hz, the coefficients are
\[
a=\frac{0.00155}{0.00155+0.0000625}
  \approx0.961240,
\]
\[
b=\frac{0.0000625}{0.00155+0.0000625}
  \approx0.038760.
\]

\topic{Audio filtering code}
\begin{lstlisting}[style=python]
from scipy.io import wavfile
import numpy as np
import matplotlib.pyplot as plt

#read sampling frequency and data from audio clip
fs, clip = wavfile.read("clip.wav")
clip = clip.astype(float); 
if clip.ndim > 1: 
    clip = clip[:,0]
clip = clip/np.max(np.abs(clip))

#RC time constant
tau = 0.00155

dt = 1.0/fs
a = tau/(tau+dt)
b = dt/(tau+dt)

y = np.zeros_like(clip)

#numeric step result convolution via DE approximation
for n in range(len(clip)): # y[n] is the current sample of the live stream
    y[n] = a*y[n-1] + b*clip[n]

#write output to files for comparison
from scipy.io import wavfile
y_i16 = np.int16(y / np.max(np.abs(y)) * 32767) # converts to 16 bit int for wav format
fname = "clip_RC_filter.wav"
wavfile.write(fname, round(fs), y_i16) # writes the wav file to disk
\end{lstlisting}

\topic{Sound compared with the FIR moving average}

The IIR filter and the FIR moving average both reduce
high-frequency content, so a similar smoother and more
muffled. Their outputs do not need to
sound exactly the same because the two filters attenuate
frequencies by different amounts.

\topic{Computational advantages and disadvantages}

The IIR update requires two multiplications and one addition
per output sample. It only needs to keep the previous output,
making it useful for real-time filtering.

Direct convolution with the 64-sample FIR moving average
requires 64 multiplications and 63 additions per output
sample. It also needs a history of input samples. This
requires more computation and memory than the first-order
IIR filter. A running-sum implementation could make this
particular moving average more efficient.

\topic{Filtering advantages and disadvantages}

The IIR filter gives an RC-like low-pass response using
only two coefficients. However, feedback means its
coefficients must be chosen carefully to keep it stable.
Here, $|a|<1$, so the filter is stable.

An FIR filter with finite coefficients is stable because
it has no feedback. Symmetric FIR filters, including
this moving average, can provide linear phase. The IIR
filter has frequency-dependent phase delay. A disadvantage
of FIR filtering is that sharper frequency selection can
require a longer filter and more computation.


\clearpage
\subsection{Task 4.3: IIR Gain versus Frequency}

The analytical RC magnitude response uses the time constant
from Task 4.2:
\[
M(\omega_0)=\frac{1}{\sqrt{1+(\omega_0\tau)^2}},
\qquad
\omega_0=2\pi f,
\qquad
\tau=0.00155\ \text{s}.
\]

\topic{Python code}
\begin{lstlisting}[style=python]
from scipy.io import wavfile
import numpy as np
import matplotlib.pyplot as plt

#sampling frequency
fs = 16000

#RC time constant
tau = 0.00155

dt = 1.0/fs
a = tau/(tau+dt)
b = dt/(tau+dt)

#frequency list
freq = [50, 100, 200, 500, 1000, 2000, 4000, 8000]

n = np.arange(0, 1.0, dt)

#FIR filter
h_ma = np.ones(64)/64 #length-64 moving average filter      

#create list of one second pure tones at each test frequency
tones = []
for f in freq:
    tones.append(np.cos(2*np.pi*f*n))

#tones filted via IIR filter derived from the first stage of the RC circuit
y_iir = []
for t in tones:
    out = np.zeros(len(t))
    for k in range(len(t)):  # k is the current sample of the live stream
        if k == 0:
            out[k] = b*t[k]
        else:
            out[k] = a*out[k-1] + b*t[k]
    y_iir.append(out)

#tones filtered via FIR filter h_ma
y_fir = []
for t in tones:
    y_fir.append(np.convolve(t, h_ma))



#measure steady state amplitude of IIR filter output at each frequency
steady_state_amp_iir = []
for out in y_iir:
    steady_state_amp_iir.append(np.max(np.abs(out[len(out)//2:])))

#measure steady state amplitude of FIR filter output at each frequency
steady_state_amp_fir = []
for out in y_fir:
    steady_state_amp_fir.append(np.max(np.abs(out[len(out)//2:])))

#analytic steady state amplitude for IIR filter derived from first stage of the RC circuit
f_analytic = np.arange(50, 8000, 1)
H_analytic = 1/np.sqrt(1 + (2*np.pi*f_analytic*tau)**2)

# code for plotting freq vs amplitude
fig, ax = plt.subplots()

ax.plot(freq, steady_state_amp_iir, 'o-', label='IIR Output Steady-State Amplitude (RC circuit)')
ax.plot(freq, steady_state_amp_fir, 'o-', label='FIR Output Steady-State Amplitude (Moving average)')
ax.plot(f_analytic, H_analytic, label='Analytic Steady-State Amplitude (RC circuit)')
ax.legend()

ax.set_title("Amplitude vs Frequency for different filter outputs")

ax.grid(True, alpha=0.3)
ax.set_xlabel("Frequency (Hz)")
ax.set_ylabel("Amplitude")
ax.set_xscale('log')


plt.tight_layout()
fig.savefig("task4.3_amplitude_vs_freq.png", dpi=200)
plt.show()
\end{lstlisting}

\labfigure{task4.3_amplitude_vs_freq.png}
{Output amplitude versus frequency for the FIR moving
average, the IIR filter, and the analytical RC response.}

\topic{Agreement between the IIR and analytical responses}

The IIR and analytical responses agree closely at the
lower test frequencies. At 100 Hz, the IIR amplitude is
about 0.709, compared with 0.716 analytically. Both curves
decrease as frequency increases.

The relative difference becomes larger near the upper end
of the frequency range. At 8000 Hz, the IIR amplitude is
about 0.0198, compared with 0.0128 analytically. The IIR
filter is a sampled approximation of the continuous-time
RC circuit, so their responses do not match exactly.
Measuring peaks only at the sample times can also affect
the measured amplitude at high frequencies.

\topic{Comparison of the three curves and low-pass behavior}

All three curves have larger amplitudes at low frequencies
and smaller amplitudes at high frequencies. Both the FIR
moving average and the IIR filter therefore show low-pass
behavior.

The two digital filters do not reduce every frequency by
the same amount. At the tested frequencies from 200 Hz
upward, the FIR output amplitude is smaller than the IIR
output amplitude. For example, at 200 Hz, the FIR amplitude
is about 0.234, while the IIR amplitude is about 0.449.
The IIR curve follows the analytical RC response more
closely because it was calculated from the RC circuit's
difference-equation approximation.

The FIR measurements include the ending transient of the
full convolution, as in Task 3.3. This contributes to the
small nonzero amplitudes at the higher test frequencies.


\section{Conclusion}

The hand calculation, convolution loop, and NumPy function
gave the same discrete-time convolution result. Convolution
with a shifted impulse delayed the signal without changing
its shape. For the RC circuit, the numerical step response
closely matched the hand calculation. The cosine response
showed that higher frequencies were reduced more than lower
frequencies. The two-stage analytical magnitudes also agreed
closely with the measurements from Lab 2.

The audio filters showed different ways to change a signal.
The moving average reduced high-frequency content, the
difference filter reduced low-frequency content, and the
echo filter added a delayed copy. Changing the order of the
moving average and echo filters gave the same result to
numerical precision. The IIR filter produced an RC-like
low-pass response with fewer calculations per sample than
direct FIR convolution. Its frequency response was close
to the analytical RC response at lower frequencies, with
larger relative differences near the sampling limit
\paragraph{Take-aways.}
\begin{itemize}
  \item Convolution gives the output of an LTI system from $x$ and $h$; the output length is $N_x+N_h-1$.
  \item Continuous-time convolution can be approximated with \texttt{np.convolve} times $dt$.
  \item An RC circuit is a low-pass filter: gain $1/\sqrt{1+(\omega\tau)^2}$, phase approaching $-90^\circ$.
  \item In a cascade, magnitudes multiply and phases add.
  \item FIR filters are always stable and can have linear phase; IIR filters are cheaper per sample but need $|a|<1$.
\end{itemize}


\end{document}